\documentclass[11pt,a4paper]{article}
\usepackage{jheppub}
\usepackage[utf8]{inputenc}
\usepackage{graphicx,fancybox,float,comment,bm}

\usepackage[dvipsnames]{xcolor}
\usepackage{amsmath,amssymb,amsfonts}
\usepackage{dsfont}
\usepackage{physics}
%\usepackage{showkeys}
\usepackage{ulem}

\title{Timelike entanglement measures for the double-scaled SYK model}

\author[a]{Johanna Erdmenger,}
\author[a]{Maximilian Gaschler,}
\author[a]{Jonathan Karl,}
\author[b]{Arnab Kundu,}
\author[a]{Georg Setzer}
\date{}
\affiliation[a]{Institute for Theoretical Physics and Astrophysics,
Julius-Maximilians-Universität Würzburg, Am Hubland, 97074 Würzburg, Germany}

\abstract{abstract}

\begin{document}
\maketitle
\newpage
\section{Introduction}



\section{A spacetime density kernel for the DSSYK model}


Given a Hilbert space $\mathcal{H}$ with density matrix $\rho$, time evolution $U_t = e^{-iHt}$,
and a bipartition $\mathcal{H} = \mathcal{H}_A \otimes \mathcal{H}_B$ with matrix unit bases
$\hat{E}^A_{ij} = \ketbra{i}{j} \otimes \mathbb{1}_B$ and
$\hat{E}^B_{kl} = \mathbb{1}_A \otimes \ketbra{k}{l} $,


\subsection{Definition and properties of the spacetime density kernel}
For two subregions in our Hilbert space we introduce the generalized spacetime density kernel GSDK $T_{AB}$ as the timelike entanglement equivalent to the normal density matrix $\rho$. The kernel encodes all time dependency of our theory and for the two-time correlators we can write
\begin{align}
	 \mathrm{Tr} \left[ T_{AB} (t) (O_{A} \otimes O_{B} ) \right] = \mathrm{Tr} \left[ \rho \ O_{A} U^{\dagger}_{t} O_{B} U_{t} \right].
	 \label{eq:GSDK_Def}
\end{align}
Furthermore it is convenient to define the components of $T_{AB}$ which we will call the response tensor
\begin{align}
	T_{AB}(t) = \sum_{i,j,k,l} C_{ij;kl} (t) E_{ji}^{A} \otimes E_{lk}^{B}.
	\label{eq:}
\end{align}
In this way the two-time correlator from \eqref{eq:GSDK_Def} becomes a simple sum over the operator components
\begin{align}
    \mathcal{L}_t(O_A, O_B)
    = \text{Tr}\left[\rho\,O_A\,U_t^\dagger\,O_B\,U_t\right]
    = \sum_{i,j,k,l} C_{ij;kl}(t)\,O^{(A)}_{ij}\,O^{(B)}_{kl},
    \label{eq:correlator_from_SDK}
\end{align}
where $O^{(A)}_{ij} = \langle i|O_A|j\rangle$ and $O^{(B)}_{kl} = \langle k|O_B|l\rangle$
are the matrix elements in the chosen basis.
We calculate the response tensor as the trace
\begin{align}
    C_{ij;kl}(t) = \text{Tr}\left[\rho\,\hat{E}^A_{ij}\,U_t^\dagger\,\hat{E}^B_{kl}\,U_t\right].
    \label{eq:SDK_def}
\end{align}

\subsection{Overview of the double-scaled SYK model}

The SYK model is consists of $N$ Majorana fermions which are coupled by a a random parameter $J$ in an $p$-body interaction.
\begin{align}
	H_{SYK} = i^{p/2} \sum_{i_1...i_p} J_{i_1 \dots i_p} \psi_{i_1} \dots \psi_{i_p}.
	\label{eq:SYK_Hamiltonian}
\end{align}
In the double scaling limit the model is exactly solvable. We take
\begin{align}
	N \rightarrow \infty  \qquad p \rightarrow \infty,
\end{align}
while keeping 
\begin{align}
	\lambda = \frac{2 p^2}{N} = \mathrm{const.}
	\label{eq:lambda}
\end{align}
constant. We also define the parameter $q = e^{-\lambda}$, which is used as the deformation parameter in the later part.\\
In this limit the moments of the Hamiltonian 
\begin{align}
    m_{k} = \ev{\mathrm{Tr}H^{k}}_{J} 
\end{align}
can be calculated as a sum over chord diagrams as showed in e.g. \cite{berkooz_towards_2018,berkooz_chord_2018}. By evaluating this sum combinatorically we go to an auxiliary Hilbert space of open chord states.
With states $\ket{l}$ as the basis of the chord number states we can define the transfer matrix 
\begin{align}
    T = 
    \begin{pmatrix}
        0 & 1 & 0 & \dots \\
        1 & 0 & \frac{}{} \\
    \end{pmatrix}
\end{align}

%The moments of the Hamiltonian are calculated by summing all the possible chord diagrams with the respective number of Hamiltonians. For each intersection of Hamiltonian chords one assigns powers of $q$ to the diagram (fig. \ref{fig:Chord_Diag_ex}).
%\begin{figure}
%%Example Chord diagram
%\label{fig:Chord_Diag_ex}
%\end{figure}
%
%
%\begin{align}
%    m_k = \sum_\text{chord diagrams} q^{\#(H \cap H)}
%    \nonumber
%\end{align}

\subsection{Derivation of the DSSYK spacetime density kernel}


\begin{align}
    C_{ij;kl}(t) = 
    \ev{0|\rho|i} \int \frac{d\theta d\phi}{(2\pi)^2} \ \mu(\theta) \mu(\phi) \
    e^{-i t(E(\theta)-E(\phi))} 
    \frac{H_j(\cos \theta |q) H_k(\cos \phi|q)}{\sqrt{(q;q)_j (q;q)_k}} \ 
    \frac{H_l(\cos\theta |q)}{\sqrt{(q;q)_l}}
\nonumber \end{align}

\subsection{Relation to the two-point function}
\section{Numerical results}
\subsection{Timelike entanglement measures}
The spacetime density kernel allows us to study correlations in time, since it encodes the dynamics of different time-slices in one operator. Multiple measures, which quantify these correlations, can be extracted from the spacetime density kernel $T_{AB}(t)$. The goal of this section is to numerically calculate these quantities, which were introduced in \cite{Das:2026ifj}, for the double-scaled SYK model. It is important to note that $T_{AB}(t)$ is not necessarily Hermitian, but $\text{Tr}\ T_{AB}(t)=1$ always holds. Firstly, we aim to calculate the Tsallis entropy
\begin{equation}
    Z_n(t)=\text{Tr} [T_{AB}(t)^n],
\end{equation}
while we mostly focus on the $n=2$ case. We also investigate the Schatten $p$-norms of $T_{AB}$, which are defined by 
\begin{equation}
    \|M\|_p=\text{Tr}\left((M^\dagger M)^{p/2}\right)^{1/p}.
\end{equation}
In particular, we are interested in the Frobenius norm 
\begin{equation}
    \|T_{AB}(t)\|_2^2=\text{Tr}\left[T_{AB}(t)T^\dagger_{AB}(t)\right].
\end{equation}
In order to quantify the non-Hermiticity of the spacetime density kernel, we also study the entanglement $p$-imagitivity 
\begin{equation}
    \mathcal{I}_p(t)=\|T_{AB}(t)-T^\dagger_{AB}(t)\|_p.
\end{equation}
Here, we focus again on the case $p=2$, since $\mathcal{I}_2(t)$ controls causal influence: it is nonzero iff signaling between A and B is possible.\\
Another important quantity which we will consider is the kernel negativity $\mathcal{N}_H(t)$, which quantifies the spectral weight of the Hermitian part $H_{AB}(t)\equiv(T_{AB}+T^\dagger_{AB})/2$ of the spacetime density kernel. It is defined by
\begin{equation}
    \mathcal{N}_H(t)=\sum_{\mu_\alpha(t)<0}|\mu_\alpha(t)|,
\end{equation}
where $\{\mu_\alpha(t)\}$ are the eigenvalues of $H_{AB}(t)$.
\subsection{Spectral form factor}
\newpage
\bibliographystyle{JHEP}
\bibliography{ref}
\end{document}